\documentclass[12pt,a4paper]{article}

% ---- Encoding & Font ----
\usepackage{iftex}
\ifPDFTeX
  \usepackage[utf8]{inputenc}
  \usepackage[T1]{fontenc}
\else
  \usepackage{fontspec}  % XeTeX/LuaTeX (tectonic): Unicode nativo (·, —, ¿, ª, §)
\fi
\usepackage[spanish]{babel}
\usepackage{csquotes}

% ---- Page layout ----
\usepackage[top=2.5cm, bottom=2.5cm, left=2.5cm, right=2.5cm]{geometry}
\usepackage{setspace}
\onehalfspacing

% ---- Math ----
\usepackage{amsmath, amssymb, amsthm}
\usepackage{mathtools}
\usepackage{physics}
\usepackage{esint}  % \oiint

% ---- Graphics ----
\usepackage{graphicx}
\usepackage{tikz}
\usepackage{float}
\usepackage{caption}

% ---- Tables ----
\usepackage{array}
\usepackage{booktabs}
\usepackage{tabularx}
\usepackage{colortbl}

% ---- Colours ----
\usepackage{xcolor}
\definecolor{notebg}{HTML}{FFF3CD}
\definecolor{noteborder}{HTML}{FFC107}
\definecolor{boxred}{HTML}{F8D7DA}
\definecolor{boxredborder}{HTML}{F5C6CB}

% ---- Boxes ----
\usepackage{mdframed}
\usepackage{tcolorbox}
\tcbuselibrary{most}

\newtcolorbox{keybox}[1][]{
  colback=blue!5,
  colframe=blue!70!black,
  arc=4pt,
  boxrule=1pt,
  fonttitle=\bfseries,
  title={#1}
}

\newtcolorbox{warnbox}[1][]{
  colback=boxred,
  colframe=boxredborder,
  coltitle=black!75,
  arc=4pt,
  boxrule=1pt,
  fonttitle=\bfseries,
  title={#1}
}

\newtcolorbox{notebox}[1][]{
  colback=notebg,
  colframe=noteborder,
  coltitle=black!75,
  arc=4pt,
  boxrule=1pt,
  fonttitle=\bfseries,
  title={#1}
}

% ---- Hyperlinks & TOC ----
\usepackage[hidelinks]{hyperref}
\usepackage{bookmark}

% ---- Enumerate ----
\usepackage{enumitem}

% ---- Miscellaneous ----
\usepackage{icomma}
\usepackage{siunitx}
\usepackage{textcomp}
\usepackage[bottom]{footmisc}

% ---- Diagramas (gráficas y esquemas) ----
\usepackage{pgfplots}
\pgfplotsset{compat=1.18}
\usetikzlibrary{babel}  % babel español activa `>`; esto evita romper las flechas `->`
% courses/ElecMag2024/Clases/assets/tikzstyles.tex
% ---------------------------------------------------------------------------
% Estilos compartidos para los diagramas del curso Electromagnetismo (2024).
% Fuente única del "look" visual (colores, grosores, escalas) para que todas
% las figuras (TikZ / pgfplots / CircuiTikz) sean consistentes entre clases.
%
% Es una COPIA de courses/Fisica3-2015/Clases/assets/tikzstyles.tex: mismo
% dominio, misma paleta. Copia y no symlink para que cada curso sea
% autocontenido y la paleta de Física III pueda evolucionar aparte.
%
% Uso: la clase debe cargar antes `tikz` y `pgfplots` (y `circuitikz` si la
%      clase tiene circuitos), y luego % courses/ElecMag2024/Clases/assets/tikzstyles.tex
% ---------------------------------------------------------------------------
% Estilos compartidos para los diagramas del curso Electromagnetismo (2024).
% Fuente única del "look" visual (colores, grosores, escalas) para que todas
% las figuras (TikZ / pgfplots / CircuiTikz) sean consistentes entre clases.
%
% Es una COPIA de courses/Fisica3-2015/Clases/assets/tikzstyles.tex: mismo
% dominio, misma paleta. Copia y no symlink para que cada curso sea
% autocontenido y la paleta de Física III pueda evolucionar aparte.
%
% Uso: la clase debe cargar antes `tikz` y `pgfplots` (y `circuitikz` si la
%      clase tiene circuitos), y luego % courses/ElecMag2024/Clases/assets/tikzstyles.tex
% ---------------------------------------------------------------------------
% Estilos compartidos para los diagramas del curso Electromagnetismo (2024).
% Fuente única del "look" visual (colores, grosores, escalas) para que todas
% las figuras (TikZ / pgfplots / CircuiTikz) sean consistentes entre clases.
%
% Es una COPIA de courses/Fisica3-2015/Clases/assets/tikzstyles.tex: mismo
% dominio, misma paleta. Copia y no symlink para que cada curso sea
% autocontenido y la paleta de Física III pueda evolucionar aparte.
%
% Uso: la clase debe cargar antes `tikz` y `pgfplots` (y `circuitikz` si la
%      clase tiene circuitos), y luego \input{../assets/tikzstyles.tex}
% (compilar desde el directorio de la clase para que la ruta relativa resuelva).
% ---------------------------------------------------------------------------

% ---- Paleta (alineada con las cajas del preámbulo: keybox/notebox/warnbox) ----
\definecolor{figblue}{HTML}{1F4E79}   % azul principal (líneas, keybox)
\definecolor{figred}{HTML}{C0392B}    % rojo de énfasis (warnbox)
\definecolor{figamber}{HTML}{B8860B}  % ámbar (notebox)
\definecolor{figgray}{HTML}{555555}   % auxiliares, ejes, guías

% ---- CircuiTikz: escala y grosor de los componentes ----
% Guarda \ifdefined: la electrostática (Clases 1-14) no carga `circuitikz`, y
% sin la guarda el \input revienta con `Undefined control sequence \ctikzset`.
% Cuando el curso llegue a circuitos, el bloque se activa solo.
\ifdefined\ctikzset
  \ctikzset{
    bipoles/thickness=1,
    resistors/scale=0.9,
    capacitors/scale=0.9,
  }
\fi

% ---- TikZ: estilos reutilizables ----
% Nota: las puntas se declaran como `-{Latex[...]}` y no como `->`. La forma
% con llaves NO contiene un `>` literal, así que es inmune al `>` activo que
% introduce babel español (ver CLAUDE.md §7.3.3).
\usetikzlibrary{arrows.meta,calc,angles,quotes}

\tikzset{
  figline/.style={figblue, line width=0.9pt},
  figaux/.style={figgray, dashed, line width=0.6pt},
  % vectores
  figvec/.style={figblue, line width=0.9pt, -{Latex[length=2.2mm,width=1.6mm]}},
  figvecres/.style={figred, line width=1.3pt, -{Latex[length=2.6mm,width=1.9mm]}},
  figvecaux/.style={figgray, line width=0.7pt, -{Latex[length=1.8mm,width=1.3mm]}},
  figaxis/.style={figgray, line width=0.6pt, -{Latex[length=2mm,width=1.4mm]}},
  % cargas puntuales
  figcarga/.style={circle, inner sep=0pt, minimum size=5.4pt, draw=none},
  figpos/.style={figcarga, fill=figred},
  figneg/.style={figcarga, fill=figblue},
  figneutra/.style={figcarga, fill=figgray},
  % etiquetas
  figlbl/.style={font=\small},
  figlblaux/.style={font=\footnotesize, figgray},
}

% ---- pgfplots: estilo base de las gráficas ----
\pgfplotsset{
  emfig/.style={
    width=11cm, height=6.5cm,
    axis lines=left,
    xlabel style={font=\small},
    ylabel style={font=\small},
    tick label style={font=\footnotesize},
    every axis plot/.append style={line width=1pt},
    grid=major,
    grid style={figgray!20},
    legend cell align=left,
    legend style={font=\footnotesize, draw=figgray!40},
  },
}

% (compilar desde el directorio de la clase para que la ruta relativa resuelva).
% ---------------------------------------------------------------------------

% ---- Paleta (alineada con las cajas del preámbulo: keybox/notebox/warnbox) ----
\definecolor{figblue}{HTML}{1F4E79}   % azul principal (líneas, keybox)
\definecolor{figred}{HTML}{C0392B}    % rojo de énfasis (warnbox)
\definecolor{figamber}{HTML}{B8860B}  % ámbar (notebox)
\definecolor{figgray}{HTML}{555555}   % auxiliares, ejes, guías

% ---- CircuiTikz: escala y grosor de los componentes ----
% Guarda \ifdefined: la electrostática (Clases 1-14) no carga `circuitikz`, y
% sin la guarda el \input revienta con `Undefined control sequence \ctikzset`.
% Cuando el curso llegue a circuitos, el bloque se activa solo.
\ifdefined\ctikzset
  \ctikzset{
    bipoles/thickness=1,
    resistors/scale=0.9,
    capacitors/scale=0.9,
  }
\fi

% ---- TikZ: estilos reutilizables ----
% Nota: las puntas se declaran como `-{Latex[...]}` y no como `->`. La forma
% con llaves NO contiene un `>` literal, así que es inmune al `>` activo que
% introduce babel español (ver CLAUDE.md §7.3.3).
\usetikzlibrary{arrows.meta,calc,angles,quotes}

\tikzset{
  figline/.style={figblue, line width=0.9pt},
  figaux/.style={figgray, dashed, line width=0.6pt},
  % vectores
  figvec/.style={figblue, line width=0.9pt, -{Latex[length=2.2mm,width=1.6mm]}},
  figvecres/.style={figred, line width=1.3pt, -{Latex[length=2.6mm,width=1.9mm]}},
  figvecaux/.style={figgray, line width=0.7pt, -{Latex[length=1.8mm,width=1.3mm]}},
  figaxis/.style={figgray, line width=0.6pt, -{Latex[length=2mm,width=1.4mm]}},
  % cargas puntuales
  figcarga/.style={circle, inner sep=0pt, minimum size=5.4pt, draw=none},
  figpos/.style={figcarga, fill=figred},
  figneg/.style={figcarga, fill=figblue},
  figneutra/.style={figcarga, fill=figgray},
  % etiquetas
  figlbl/.style={font=\small},
  figlblaux/.style={font=\footnotesize, figgray},
}

% ---- pgfplots: estilo base de las gráficas ----
\pgfplotsset{
  emfig/.style={
    width=11cm, height=6.5cm,
    axis lines=left,
    xlabel style={font=\small},
    ylabel style={font=\small},
    tick label style={font=\footnotesize},
    every axis plot/.append style={line width=1pt},
    grid=major,
    grid style={figgray!20},
    legend cell align=left,
    legend style={font=\footnotesize, draw=figgray!40},
  },
}

% (compilar desde el directorio de la clase para que la ruta relativa resuelva).
% ---------------------------------------------------------------------------

% ---- Paleta (alineada con las cajas del preámbulo: keybox/notebox/warnbox) ----
\definecolor{figblue}{HTML}{1F4E79}   % azul principal (líneas, keybox)
\definecolor{figred}{HTML}{C0392B}    % rojo de énfasis (warnbox)
\definecolor{figamber}{HTML}{B8860B}  % ámbar (notebox)
\definecolor{figgray}{HTML}{555555}   % auxiliares, ejes, guías

% ---- CircuiTikz: escala y grosor de los componentes ----
% Guarda \ifdefined: la electrostática (Clases 1-14) no carga `circuitikz`, y
% sin la guarda el \input revienta con `Undefined control sequence \ctikzset`.
% Cuando el curso llegue a circuitos, el bloque se activa solo.
\ifdefined\ctikzset
  \ctikzset{
    bipoles/thickness=1,
    resistors/scale=0.9,
    capacitors/scale=0.9,
  }
\fi

% ---- TikZ: estilos reutilizables ----
% Nota: las puntas se declaran como `-{Latex[...]}` y no como `->`. La forma
% con llaves NO contiene un `>` literal, así que es inmune al `>` activo que
% introduce babel español (ver CLAUDE.md §7.3.3).
\usetikzlibrary{arrows.meta,calc,angles,quotes}

\tikzset{
  figline/.style={figblue, line width=0.9pt},
  figaux/.style={figgray, dashed, line width=0.6pt},
  % vectores
  figvec/.style={figblue, line width=0.9pt, -{Latex[length=2.2mm,width=1.6mm]}},
  figvecres/.style={figred, line width=1.3pt, -{Latex[length=2.6mm,width=1.9mm]}},
  figvecaux/.style={figgray, line width=0.7pt, -{Latex[length=1.8mm,width=1.3mm]}},
  figaxis/.style={figgray, line width=0.6pt, -{Latex[length=2mm,width=1.4mm]}},
  % cargas puntuales
  figcarga/.style={circle, inner sep=0pt, minimum size=5.4pt, draw=none},
  figpos/.style={figcarga, fill=figred},
  figneg/.style={figcarga, fill=figblue},
  figneutra/.style={figcarga, fill=figgray},
  % etiquetas
  figlbl/.style={font=\small},
  figlblaux/.style={font=\footnotesize, figgray},
}

% ---- pgfplots: estilo base de las gráficas ----
\pgfplotsset{
  emfig/.style={
    width=11cm, height=6.5cm,
    axis lines=left,
    xlabel style={font=\small},
    ylabel style={font=\small},
    tick label style={font=\footnotesize},
    every axis plot/.append style={line width=1pt},
    grid=major,
    grid style={figgray!20},
    legend cell align=left,
    legend style={font=\footnotesize, draw=figgray!40},
  },
}


% ---- Title ----
\title{\textbf{Clase 21: Propiedades magnéticas de la materia (parte 2)} \\[0.3em]
        \large{Intensidad magnética, materiales lineales, ferromagnetismo, histéresis y condiciones de borde}}
\author{Transcripto y expandido --- Curso Electromagnetismo (OpenFING)}
\date{}

\begin{document}

\maketitle
\thispagestyle{empty}
\newpage

\tableofcontents
\newpage

\section{Repaso: corrientes de magnetización}

La clase anterior empezó el estudio de las propiedades magnéticas de la
materia. Es todavía más importante que en el caso eléctrico, porque el
magnetismo conocido desde la antigüedad es el de los imanes, no el de las
corrientes.

El resultado principal: el potencial vector producido por una magnetización
$\vec M$ (momento dipolar magnético por unidad de volumen) se escribe como el
de una corriente de conducción, vía Biot y Savart, sólo que con \textbf{corrientes de
magnetización}: una densidad volumétrica en el interior del material,

\[
\vec J_M = \curl\vec M,
\]
y típicamente otra superficial en el borde.

Ahora se trata de escribir las ecuaciones del campo en presencia del material
de modo que esas corrientes ---a las que no se tiene acceso experimental--- no
aparezcan, o aparezcan lo menos posible. El docente insiste en la imagen de
\textbf{barrer la mugre bajo la alfombra}: la mugre sigue ahí; la pregunta es en qué
medida se puede evitar conocerla en detalle.

\section{Ecuaciones de campo en un material}

\subsection{Campo microscópico y campo macroscópico}

Dentro de un material no interesa el campo \textbf{microscópico}, medido a escala
atómica, sino el \textbf{macroscópico}. Sólo para éste se dedujo la forma del
potencial vector dentro del material: la deducción reemplazaba todas las
corrientes por dipolos, lo que vale lejos de ellas o promediando en volúmenes
grandes.

El campo, microscópico o macroscópico, verifica siempre

\[
\div\vec B = 0, \qquad \curl\vec B = \mu_0\,\vec J_{\text{total}},
\]
donde en la ley de Ampère entran \textbf{todas} las corrientes: las de conducción y
las microscópicas, asociadas a los espines o a las corrientes atómicas:

\[
\vec J_{\text{total}} = \vec J + \vec J_M,
\]
con $\vec J$ la corriente de conducción. Esto vale en el interior del material;
las corrientes de magnetización superficiales del borde se tratan aparte.

La clave para el campo macroscópico es que se conoce la forma explícita de
$\vec J_M$ en términos de la magnetización:

\[
\curl\vec B = \mu_0\,\vec J + \mu_0\,\curl\vec M .
\]

\subsection{La intensidad magnética H}

Dividiendo por $\mu_0$ y pasando el rotacional al otro lado:

\[
\curl\left(\frac{\vec B}{\mu_0} - \vec M\right) = \vec J .
\]

Se define la \textbf{intensidad magnética}

\begin{keybox}
\[
\boxed{\ \vec H = \frac{\vec B}{\mu_0} - \vec M\ }
\]
\end{keybox}
y queda

\begin{keybox}
\[
\boxed{\ \curl\vec H = \vec J\ }
\]
\end{keybox}

Es como la ley de Ampère, pero para $\vec H$ y \textbf{sólo con la corriente de
conducción}, que es la que realmente se controla. Es la \textbf{ley de Ampère en un
material} (en su interior; las corrientes superficiales se tratan aparte).

Es exactamente lo que se hizo en el caso eléctrico al pasar de $\vec E$ a
$\vec D$ metiendo adentro la polarización. «Por ahora no estamos haciendo más
que esconder la mugre, poniéndole un nombre.»

\textbf{Unidades.} $\vec H$ y $\vec M$ tienen las mismas unidades. Como
$\curl\vec H = \vec J$ y $\vec J$ se mide en A/m$^2$,

\[
[H] = [M] = \text{A/m}.
\]

\begin{notebox}
El docente se equivoca primero con las unidades («metro cúbico» en algún
lado) y lo corrige media hora después: son amperes por metro.
\end{notebox}

\subsection{Lo que falta: la relación constitutiva}

«Hasta ahora esto es lindo e inútil.» Hay una ecuación para $\vec B$
(divergencia nula) y otra para $\vec H$ (Ampère); sin una relación entre los
dos campos el problema no se resuelve. Igual que en el caso eléctrico había
que saber la relación entre $\vec E$ y $\vec D$, acá alguien tiene que dar

\[
\vec B = \vec B(\vec H) \quad\text{o}\quad \vec H = \vec H(\vec B):
\]
la \textbf{relación constitutiva} o \textbf{ecuación de estado} del material.

\begin{itemize}
  \item \textbf{Idealmente} se deduce de la física microscópica: un número del orden del
  de Avogadro de espines y momentos orbitales interactuando entre sí y con el
  campo externo, tratados con \textbf{mecánica estadística}. En general es un
  problema muy difícil; se sabe resolver cuando las interacciones entre átomos
  son despreciables o suficientemente simples. Es toda una rama de la física y
  no se hace en el curso.
  \item \textbf{En la práctica} se toma de datos experimentales. Para un material
  homogéneo e isotrópico $\vec B$ es paralelo a $\vec H$ (no hay direcciones
  preferenciales) y alcanza una \textbf{tabla} ---una curva--- que dé el módulo de
  $\vec B$ para cada módulo de $\vec H$.
\end{itemize}

\subsection{Forma integral de la ley de Ampère}

Se toma una curva cerrada $C$ (una curva amperiana) y una superficie $S$ que
se apoya en ella, orientada con la regla de la mano derecha. Por el teorema de
Stokes y la ley de Ampère en el material:

\[
\oint_C \vec H\cdot d\vec l = \int_S (\curl\vec H)\cdot\hat n\,dS
= \int_S \vec J\cdot\hat n\,dS
\]

\begin{keybox}
\[
\boxed{\ \oint_C \vec H\cdot d\vec l = I_{\text{adentro}}\ }
\]
\end{keybox}

\begin{warnbox}
\textbf{Sólo cuenta la corriente de conducción} que atraviesa $S$; las corrientes
de magnetización no.
\end{warnbox}

La otra ecuación en forma integral ya se conoce: el flujo de $\vec B$ por
cualquier superficie cerrada es cero.

\section{Susceptibilidad y permeabilidad: materiales lineales}

\subsection{Definición}

Primero el caso más simple, parecido al eléctrico. Se suponen materiales
\textbf{homogéneos e isotrópicos}: $\vec B$, $\vec M$ y $\vec H$ son paralelos, hay
una sola dirección en juego. Muchos metales tienen una red cristalina con
direcciones preferenciales, pero en la práctica son policristalinos ---muchos
cristalitos orientados al azar--- y macroscópicamente se comportan como
isotrópicos.

Un material es \textbf{lineal} si

\begin{keybox}
\[
\boxed{\ \vec M = \chi_m\,\vec H\ }
\]
\end{keybox}
con $\chi_m$ la \textbf{susceptibilidad magnética}. Como $\vec M$ y $\vec H$ tienen
las mismas unidades, $\chi_m$ es \textbf{adimensionada}: se compara con 1.

\begin{notebox}
El docente arranca escribiendo $\vec M = \chi_m\vec B$ y se corrige: hay dos
definiciones en la literatura, y la que usa el curso es con $\vec H$.
\end{notebox}

El nombre es «de lo más apropiado»: una persona susceptible se ofende mucho
por poca cosa; un material con $\chi_m$ grande reacciona mucho a un campo
chico.

\begin{notebox}
\textbf{Si no es isotrópico}, en el caso lineal y homogéneo la susceptibilidad es
una \textbf{matriz} (un tensor): se aplica el campo en una dirección y la
magnetización aparece en otra. Si además no es homogéneo, la matriz depende
del punto.
\end{notebox}

\subsection{Diamagnetismo y paramagnetismo}

Diferencia fundamental con el caso eléctrico, donde la susceptibilidad era
siempre positiva (la polarización va para el mismo lado que el campo):

\begin{center}
\renewcommand{\arraystretch}{1.35}
\small
\begin{tabularx}{\textwidth}{l>{\raggedright\arraybackslash}X>{\raggedright\arraybackslash}X}
\toprule
\textbf{$\chi_m$} & \textbf{Nombre} & \textbf{Qué pasa} \\
\midrule
$< 0$ & \textbf{Diamagnetismo} & $\vec M$ apunta en sentido opuesto a $\vec H$ \\
$> 0$ & \textbf{Paramagnetismo} & $\vec M$ apunta en el mismo sentido que $\vec H$ \\
\bottomrule
\end{tabularx}
\end{center}

Típicamente la reacción de los \textbf{espines} es paramagnética y la de las
\textbf{corrientes orbitales} es diamagnética. Suelen estar las dos y normalmente
gana la paramagnética; pero en materiales con tantos electrones de un espín
como del opuesto el efecto paramagnético se anula y sobrevive el diamagnético.
Quién gana depende de la estructura molecular.

\begin{notebox}
$\chi_m$ es constante \textbf{a una temperatura dada}, pero en general depende
bastante de ella: a temperaturas muy altas tiende a cero, y en la mayoría de
los materiales magnéticos hay una temperatura ---la \textbf{temperatura de Curie}---
en la que diverge. El curso no trata la relación entre propiedades térmicas y
magnetismo.
\end{notebox}

\subsection{Órdenes de magnitud}

A la pregunta de un estudiante sobre el valor de $\chi_m$, el docente responde
que a temperatura ambiente «básicamente no existe»: existe, pero es
chiquitito. Valores dados en clase:

\begin{center}
\renewcommand{\arraystretch}{1.35}
\small
\begin{tabularx}{\textwidth}{l>{\raggedright\arraybackslash}X>{\raggedright\arraybackslash}X}
\toprule
\textbf{Material} & \textbf{$\chi_m$} & \textbf{Tipo} \\
\midrule
Aluminio & $2{,}1\times10^{-5}$ & para \\
Bismuto & $-16{,}4\times10^{-5}$ & dia \\
Cobre & $-0{,}98\times10^{-5}$ & dia \\
Diamante & $-2{,}2\times10^{-5}$ & dia \\
Magnesio & $1{,}2\times10^{-5}$ & para \\
Oxígeno (1 atm) & $0{,}193\times10^{-5}$ & para \\
Titanio & $18{,}0\times10^{-5}$ & para \\
Tungsteno & $7{,}6\times10^{-5}$ & para \\
\bottomrule
\end{tabularx}
\end{center}

Van de $10^{-5}$ a $10^{-4}$. Como

\[
\vec B = \mu_0\left(\vec H + \vec M\right),
\]
y $\vec M$ es $10^{-5}$ o $10^{-4}$ veces $\vec H$, estos materiales se
comportan magnéticamente \textbf{casi como el vacío}.

\begin{notebox}[Los lineales son aburridos]
Existen, son comunes, la aproximación lineal
es muy buena, pero la respuesta es «chirolas, mucho menos que una propina».
Muy distinto del caso eléctrico, donde los dieléctricos lineales tienen
constantes dieléctricas de 2, 3 o 10.
\end{notebox}

\subsection{Permeabilidad magnética}

Si $\vec M$ es lineal en $\vec H$, también lo es $\vec B$, porque la relación
de arriba es exacta:

\[
\vec B = \mu_0\,(1+\chi_m)\,\vec H
\quad\Longrightarrow\quad
\boxed{\ \vec B = \mu\,\vec H,\qquad \mu = \mu_0\,(1+\chi_m)\ }
\]

$\mu$ es la \textbf{permeabilidad magnética} del material. Generaliza $\mu_0$: sin
material, $\mu = \mu_0$.

En resumen: se puede barrer la mugre pasando de $\vec B$ a $\vec H$, a costa
de necesitar la relación entre ellos. En el caso lineal todo queda en una
constante y se resuelve como en el caso eléctrico; lamentablemente el efecto
es chiquitito.

\section{Ferromagnetismo}

Los materiales con efectos magnéticos \textbf{grandes} son, casi siempre, \textbf{no
lineales}. El caso más común es el \textbf{ferromagnetismo}, cuyo nombre viene del
hierro: «el magnetismo del tipo del hierro».

\subsection{De dónde sale}

El cuento en pocas palabras, según el docente:

\begin{enumerate}
  \item El ferromagnetismo está dominado por los \textbf{espines} de los electrones.
  Normalmente los orbitales atómicos se llenan con un electrón «para arriba» y
  otro «para abajo», y sus efectos magnéticos se compensan.
  \item En el cristal de hierro hay un electrón \textbf{desapareado}, no de conducción,
  bastante ligado a la red. Su momento dipolar tiende a alinearse con un campo
  externo. Hasta ahí es como el paramagnetismo, un efecto básicamente lineal.
  \item Lo no trivial: por un \textbf{efecto cuántico complicado} ---no la interacción
  dipolo-dipolo--- los espines de átomos vecinos tienden a \textbf{alinearse entre
  sí}. Energéticamente es favorable que estén alineados.
  \item Resultado: la enorme mayoría de los espines del material se alinean unos con
  otros, y aparece una magnetización grande \textbf{aun sin campo externo}. Por eso
  hay imanes.
\end{enumerate}

\begin{notebox}
Hay otras formas de magnetismo no lineal. En el \textbf{antiferromagnetismo} la
interacción entre espines vecinos tiende a \textbf{desalinearlos}. Existe
bastante en la naturaleza, en compuestos más complicados; no se trata en el
curso.
\end{notebox}

\begin{notebox}
El docente intercala una digresión sin cerrar: macroscópicamente uno espera
que todo movimiento se degrade en calor (rozamiento), pero a nivel molecular,
donde manda la mecánica cuántica, puede haber corrientes que no se detienen.
\end{notebox}

\subsection{Características}

Para el caso homogéneo e isotrópico (realista si el material es
policristalino):

\begin{enumerate}
  \item $\vec M$ es paralela a $\vec B$ y \textbf{del mismo sentido}.
  \item Si se reduce $\vec B$ a cero, $\vec M$ \textbf{no tiende a cero}: hay
  magnetización espontánea.
  \item El efecto es \textbf{enorme}.
\end{enumerate}

Aunque la relación no sea lineal, se puede definir una \textbf{función}
permeabilidad:

\[
\vec B = \mu(H)\,\vec H,
\]
que ahora depende del módulo de $\vec H$; por eso la relación no es lineal.

\begin{warnbox}[Esta definición es problemática]
Supone que la relación es entre el
$\vec B$ y el $\vec H$ de aquí y ahora. En los ferromagnéticos depende
además de la \textbf{historia} del material (§5).
\end{warnbox}

\begin{center}
\renewcommand{\arraystretch}{1.35}
\small
\begin{tabularx}{\textwidth}{>{\raggedright\arraybackslash}X>{\raggedright\arraybackslash}X}
\toprule
\textbf{Material} & \textbf{$\mu$} \\
\midrule
Diamagnético o paramagnético & $\mu\approx\mu_0$ \\
Ferromagnético & $\mu/\mu_0$ varía mucho con $H$ y suele alcanzar $10^4$ o $10^5$ \\
\bottomrule
\end{tabularx}
\end{center}

\subsection{Curva de primera imantación}

El docente dibuja, para \textbf{hierro inicialmente desmagnetizado}, $B$ (en tesla,
«una unidad enorme») en función de $H$ (en A/m, con marcas en 200 y 400 A/m):

\begin{itemize}
  \item arranca más o menos lineal;
  \item después crece más rápido;
  \item después empieza a \textbf{saturar}, con una pendiente mucho más baja; la zona de
  saturación está del orden de \textbf{un tesla}, un campo muy grande.
\end{itemize}

\begin{figure}[H]
  \centering
  % Clase 21 — curva de primera imantación del hierro inicialmente desmagnetizado,
% como la dibuja el docente: B(H) arranca casi lineal, crece más rápido y
% satura (del orden de 1 T); mu/mu0 tiene un máximo (~6000) cerca de 200 A/m y
% luego decrece. CUALITATIVA: las curvas son funciones de forma, no datos.
% B = 1.75 tanh((H/180)^1.6) T, y mu/mu0 = B/(mu0 H), con máximo ~5900.
\begin{tikzpicture}
  \begin{axis}[
    emfig, width=6.4cm, height=5cm,
    xmin=0, xmax=600, ymin=0, ymax=2.1,
    xtick={0,200,400}, ytick={0,0.5,1,1.5,2},
    xlabel={$H$ (A/m)}, ylabel={$B$ (T)},
    title={\small (a) $B(H)$},
  ]
    \addplot[figblue, domain=0:600, samples=120] {1.75*tanh((x/180)^1.6)};
  \end{axis}
  \begin{axis}[
    emfig, width=6.4cm, height=5cm, xshift=7.2cm,
    xmin=0, xmax=600, ymin=0, ymax=7000,
    xtick={0,200,400}, ytick={0,2000,4000,6000},
    scaled y ticks=false,
    yticklabel style={/pgf/number format/1000 sep={\,}},
    xlabel={$H$ (A/m)}, ylabel={$\mu/\mu_0$},
    title={\small (b) $\mu(H)/\mu_0$},
  ]
    \addplot[figred, domain=1:600, samples=150]
      {1.75*tanh((x/180)^1.6)/(4*pi*1e-7*x)};
  \end{axis}
\end{tikzpicture}

  \caption{Curva de primera imantación del hierro, cualitativa, como la dibuja el docente: (a) $B$ satura del orden de un tesla; (b) $\mu/\mu_0$ tiene un máximo de algunos miles cerca de 200 A/m y luego decrece hacia 1. Las curvas son funciones de forma, no datos medidos.}
\end{figure}

La curva de $\mu/\mu_0$ en función de $H$ arranca en un valor moderado, tiene un
\textbf{máximo} alrededor de 200 A/m ---«muy alto, tipo 6000»--- y luego decrece hacia
1. Propiedades:

\begin{enumerate}
  \item $\mu$ es \textbf{positiva} en esta curva;
  \item \textbf{varía mucho} con $H$;
  \item \textbf{tiende a $\mu_0$} para $H$ grande.
\end{enumerate}

\textbf{Por qué tiende a $\mu_0$.} Con un campo muy grande todos los espines quedan
alineados: la magnetización \textbf{satura} y no puede crecer más, $\vec M\to$
constante. Entonces, cuando $H\gg M$,

\[
\vec B = \mu_0(\vec H+\vec M)\approx\mu_0\,\vec H,
\]
como si no hubiera magnetización. Pero hay que llegar a campos gigantescos,
difíciles de lograr en un laboratorio; la situación normal está antes de la
saturación.

\begin{notebox}
El docente aclara que su dibujo exagera: comparado con 6000, el valor final
1 es casi cero.
\end{notebox}

\subsection{Magnetización de saturación}

Es la magnetización máxima que se puede alcanzar en un material. Valores dados
en clase, como $\mu_0 M_s$ (que tiene unidades de campo magnético):

\begin{center}
\renewcommand{\arraystretch}{1.35}
\small
\begin{tabularx}{\textwidth}{>{\raggedright\arraybackslash}X>{\raggedright\arraybackslash}X}
\toprule
\textbf{Material} & \textbf{$\mu_0 M_s$ (T)} \\
\midrule
Hierro & 2,15 \\
Cobalto & 1,79 \\
Níquel & 0,61 \\
Ferrita de manganeso, MnFe$_2$O$_4$ & 0,49 \\
Ferrita de níquel, NiFe$_2$O$_4$ & 0,32 \\
\bottomrule
\end{tabularx}
\end{center}

\begin{notebox}
La transcripción nombra los dos últimos como «magnetita de magnesio» y
«magnetita de níquel»; el docente dicta sus fórmulas (manganeso, hierro dos,
oxígeno cuatro; níquel, hierro dos, oxígeno cuatro), que son las ferritas de
la tabla.
\end{notebox}

Son los ferromagnéticos más comunes: hierro, cobalto, níquel y óxidos de
hierro. En la naturaleza lo frecuente ---y lo que se conocía en la antigüedad---
son los óxidos, porque el hierro se oxida muy fácil.

\section{Histéresis}

\subsection{El material recuerda su historia}

Si la relación dependiera sólo del presente, al bajar el campo después de
subirlo se volvería por la misma curva y, con $H = 0$, se tendría $B = 0$ y
$M = 0$: el material quedaría desmagnetizado. Pero estos materiales se
\textbf{imantan}. Ese fenómeno ---depender de la historia pasada--- se llama
\textbf{histéresis}.

\subsection{Ciclos de histéresis}

Partiendo de la curva de primera imantación hasta cierto punto y bajando $H$:

\begin{itemize}
  \item el sistema \textbf{no vuelve por la curva original}, sino por encima;
  \item con $H = 0$ (sin corrientes de conducción externas) queda un campo $B$ no
  nulo adentro: es $\mu_0$ por la \textbf{magnetización espontánea};
  \item hay un tramo en que $\mu$ es \textbf{negativa} (en la curva de primera
  imantación era positiva);
  \item aplicando $H$ hacia el otro lado, en algún momento la magnetización se da
  vuelta y la curva sigue; volviendo a subir se recorre otra rama, y se cierra
  un \textbf{ciclo}.
\end{itemize}

\begin{center}
\renewcommand{\arraystretch}{1.35}
\small
\begin{tabularx}{\textwidth}{l>{\raggedright\arraybackslash}X}
\toprule
\textbf{Punto} & \textbf{Nombre} \\
\midrule
$B$ con $H = 0$ & \textbf{Remanencia} \\
$H$ con $B = 0$ & \textbf{Coercitividad}: cuánto campo externo hay que aplicar hacia el otro lado para anular el efecto neto \\
\bottomrule
\end{tabularx}
\end{center}

\begin{warnbox}[No hay una única curva]
Hay muchos ciclos de histéresis, según el campo
máximo que se llegó a aplicar; en cuál se está depende de la historia. Sólo
si se llega a \textbf{saturar} el material el ciclo es único. Además los ciclos
deberían ser simétricos (la curva es impar); el docente aclara que su dibujo
no lo es.
\end{warnbox}

Los ciclos tienen propiedades \textbf{disipativas}, con consecuencias
termodinámicas importantes que no se tratan.

La dificultad es doble: la relación no es lineal \textbf{y} no es única.

\begin{figure}[H]
  \centering
  % Clase 21 — ciclo de histéresis saturado (cualitativo) y curva de primera
% imantación punteada. Remanencia: B con H = 0. Coercitividad: H con B = 0.
% El ciclo es simétrico (curva impar), a diferencia del dibujo del pizarrón.
% Ramas: B = Bs tanh((H -+ Hc)/w), con Bs = 1.3, Hc = 1, w = 0.9.
\begin{tikzpicture}
  \begin{axis}[
    emfig, width=9.5cm, height=7cm,
    axis lines=middle,
    xlabel style={anchor=north west, font=\small},
    ylabel style={anchor=south east, font=\small},
    xmin=-4.2, xmax=4.2, ymin=-1.65, ymax=1.65,
    xtick=\empty, ytick=\empty, grid=none, clip=false,
    xlabel={$H$}, ylabel={$B$},
  ]
    \addplot[figblue, domain=-4:4, samples=120] {1.3*tanh((x+1)/0.9)};
    \addplot[figblue, domain=-4:4, samples=120] {1.3*tanh((x-1)/0.9)};
    \addplot[figgray, dashed, domain=0:3.2, samples=60] {1.3*tanh((x/1.6)^1.5)};
    % flechas de recorrido
    \draw[figvec] (axis cs:2.2,1.28) -- (axis cs:1.2,1.18);
    \draw[figvec] (axis cs:-2.2,-1.28) -- (axis cs:-1.2,-1.18);
    % remanencia y coercitividad
    \fill[figred] (axis cs:0,{1.3*tanh(1/0.9)}) circle (2pt);
    \draw[figred, line width=0.5pt] (axis cs:0,1.04) -- (axis cs:-1.1,1.4);
    \node[font=\footnotesize, text=figred, anchor=east] at (axis cs:-1.1,1.4) {remanencia};
    \fill[figred] (axis cs:-1,0) circle (2pt);
    \draw[figred, line width=0.5pt] (axis cs:-1,0) -- (axis cs:-2.0,-0.45);
    \node[font=\footnotesize, text=figred, anchor=north east] at (axis cs:-2.0,-0.45) {coercitividad};
    \draw[figgray, line width=0.5pt] (axis cs:1.9,0.95) -- (axis cs:2.5,0.45);
    \node[font=\footnotesize, text=figgray, anchor=north west] at (axis cs:2.45,0.45) {primera imantación};
  \end{axis}
\end{tikzpicture}

  \caption{Ciclo de histéresis saturado (cualitativo) y, punteada, la curva de primera imantación. En el pizarrón el ciclo no era simétrico; acá se lo dibuja impar, como aclara el docente que debería ser.}
\end{figure}

\subsection{Cómo se hacen cuentas en la práctica}

Tres situaciones en las que se puede calcular:

\begin{center}
\renewcommand{\arraystretch}{1.35}
\small
\begin{tabularx}{\textwidth}{>{\raggedright\arraybackslash}X>{\raggedright\arraybackslash}X>{\raggedright\arraybackslash}X}
\toprule
\textbf{Situación} & \textbf{Aproximación} & \textbf{Comentario} \\
\midrule
Ciclos chicos, cerca del origen & Lineal, $\vec B\approx\mu\vec H$ con $\mu$ grande (del orden de mil) & Es «una recta gorda»: aproximación de calidad regular, pero con efecto magnético genuino y cuentas simples \\
Se conoce la curva & Se reemplaza el ciclo por tramos lineales, p. ej. un \textbf{paralelogramo} & Hay que dar la curva, no un número; aparece en el práctico \\
Magnetización saturada, campo aplicado chico & $\vec M\approx$ constante $=\vec M_s$: $\vec B = \mu_0(\vec H+\vec M_s)$ & Relación \textbf{afín}: no es proporcional, pero depende sólo de una constante \\
\bottomrule
\end{tabularx}
\end{center}

\begin{notebox}
En la relación afín el docente escribe primero $\mu$ y corrige a $\mu_0$.
Aclara también que vale componente a componente: el módulo no se puede sumar
así porque el signo puede cambiar.
\end{notebox}

La utilidad tecnológica es clara: para efectos magnéticos grandes hace falta un
material ferromagnético. Por ejemplo, los \textbf{transformadores} eficientes
llevan un núcleo de hierro, y para estudiarlos habrá que ocuparse de esta
relación.

\begin{notebox}
Al cerrar el tema el docente dice que los materiales «diamagnéticos y
ferromagnéticos» son, en primera aproximación, como aire; por el contexto se
refiere a los \textbf{diamagnéticos y paramagnéticos}.
\end{notebox}

\section{Condiciones de borde}

Sección 7.6, parecida al caso eléctrico. Dos medios, 1 y 2 (uno puede ser el
vacío). Se usan las dos construcciones de siempre: una \textbf{superficie gaussiana}
$S$ con forma de cajita chata y una \textbf{curva amperiana} $C$ rectangular, ambas
a caballo de la interfaz.

\begin{figure}[H]
  \centering
  % Clase 21 — condiciones de borde magnéticas. Cajita chata S a caballo de la
% interfaz (flujo de B nulo => B_n continua) y curva amperiana C con lado L0
% en el medio 1 (Ampère en el material => salto de H_t igual a la corriente
% superficial de conducción j, que acá sale del papel). n apunta de 2 a 1.
\begin{tikzpicture}[scale=1]

  \fill[figgray!12] (-4.2,0) rectangle (4.2,1.9);
  \fill[figgray!28] (-4.2,-1.9) rectangle (4.2,0);
  \draw[figline] (-4.2,0) -- (4.2,0);
  \node[figlbl] at (-3.65,1.5) {$1$};
  \node[figlbl] at (-3.65,-1.5) {$2$};

  % normal
  \draw[figvecaux] (-3.0,0) -- (-3.0,0.85);
  \node[figlblaux] at (-2.7,0.75) {$\hat n$};

  % cajita S
  \draw[figamber, line width=0.9pt, fill=figamber!15] (-1.9,-0.16) rectangle (-0.3,0.16);
  \node[figlbl, text=figamber] at (-0.05,0.35) {$S$};
  \node[figlbl] at (-1.1,1.0) {$B_{1n} = B_{2n}$};

  % curva C
  \draw[figred, line width=0.9pt] (1.0,-0.32) rectangle (3.2,0.32);
  \draw[figred, line width=0.9pt, -{Latex[length=1.8mm,width=1.3mm]}] (1.8,0.32) -- (2.4,0.32);
  \draw[figred, line width=0.9pt, -{Latex[length=1.8mm,width=1.3mm]}] (2.4,-0.32) -- (1.8,-0.32);
  \node[figlbl, text=figred] at (2.1,0.62) {$\vec L_0$};
  \node[figlbl, text=figred] at (3.5,0.3) {$C$};

  % corriente superficial saliente
  \foreach \jx in {1.45,2.1,2.75} {
    \draw[figblue, line width=0.8pt] (\jx,0) circle (0.09);
    \fill[figblue] (\jx,0) circle (0.03);
  }
  \node[figlbl] at (2.1,1.25) {$\vec H_{1t} - \vec H_{2t} = \vec j\times\hat n$};
  \node[figlblaux] at (2.1,-0.7) {$\vec j$ (sale del papel)};

\end{tikzpicture}

  \caption{Condiciones de borde: el flujo de $\vec B$ por la cajita chata $S$ da la continuidad de $B_n$; la circulación de $\vec H$ en la curva $C$, el salto de $\vec H_t$ por una corriente superficial de conducción $\vec j$. La normal $\hat n$ apunta del medio 2 al 1.}
\end{figure}

\subsection{Componente normal de B}

El flujo de $\vec B$ por la cajita es cero. Haciéndola chata:

\begin{keybox}
\[
\boxed{\ B_{1n} = B_{2n}\ }
\]
\end{keybox}

En el caso eléctrico aparecía además la densidad superficial de carga; acá no
hay nada porque \textbf{no existen cargas magnéticas}.

\subsection{Componente tangencial de H}

Se aplica Ampère en el material a la curva $C$, con lado $\vec L_0$ en el medio
1 (y $-\vec L_0$ en el 2) y lados cortos despreciables:

\[
\oint_C \vec H\cdot d\vec l = \vec L_0\cdot(\vec H_1-\vec H_2) = I_{\text{adentro}} .
\]

Normalmente la corriente encerrada tiende a cero al achicar la curva. Pero
puede haber una \textbf{corriente de conducción concentrada en la superficie}, con
densidad superficial $\vec j$ (por ejemplo, si el medio es un conductor con
corriente circulando por su superficie). El docente lo incluye para hacer el
caso general, aunque «muchas veces suele ser cero».

La corriente que cruza la curva es la componente de $\vec j$ perpendicular al
plano de $C$, por el largo $L_0$. Con un vector normal a la interfaz
$\hat n$, el vector $\hat n\times\vec L_0$ es perpendicular al plano de la
curva y de módulo $L_0$; usando la propiedad cíclica del producto mixto,

\[
\vec L_0\cdot(\vec H_1-\vec H_2) = \vec j\cdot(\hat n\times\vec L_0)
= \vec L_0\cdot(\vec j\times\hat n).
\]

Como vale para \textbf{cualquier} $\vec L_0$ en el plano tangente:

\begin{keybox}
\[
\boxed{\ \vec H_{1t} - \vec H_{2t} = \vec j\times\hat n\ }
\]
\end{keybox}
donde $\vec j\times\hat n$ ya es tangente, sin necesidad de proyectarlo.

\begin{notebox}[El sentido de $\hat n$]
El docente redibuja la figura varias veces y
cambia el sentido de la normal sobre la marcha; termina la clase sin fijarlo
por escrito. Con la orientación de la curva de arriba, para que el resultado
sea correcto $\hat n$ tiene que apuntar \textbf{del medio 2 al medio 1}.
\end{notebox}

Quedan, como siempre, una condición normal y una tangencial. Sin corriente
superficial de conducción ($\vec j = 0$), la componente tangencial de $\vec H$
es continua.

\noindent\emph{La próxima clase sigue el martes.}

\end{document}
